\chapter{Esophageal Manometry}

\section*{Overview}
Esophageal manometry measures the strength and coordination of muscle contractions in the esophagus and the function of its upper and lower sphincters. The exact approach, setting, and preparation depend on the indication, anatomy, and the patient's health.

\section*{Why It Is Done}
It is used to evaluate difficulty swallowing, chest pain not caused by the heart, regurgitation, or suspected motility disorders such as achalasia. It may also be performed before selected anti-reflux procedures.

\section*{How It Is Performed}
After the nose and throat are numbed, a thin pressure-sensing catheter is passed through the nose into the stomach. The patient swallows small amounts of water while the catheter records pressure and muscle coordination. The catheter is then removed.

\section*{Preparation}
The patient usually fasts for several hours. Some medicines that affect esophageal movement may need to be paused under medical guidance. Tell the clinician about nasal obstruction, bleeding problems, or medicines that increase bleeding risk.

\section*{Risks and Results}
Temporary nasal or throat discomfort, gagging, watering of the eyes, and minor nosebleeding can occur. Results show the timing and strength of contractions and sphincter relaxation. Abnormal patterns may support diagnoses such as achalasia or esophageal spasm.

\section*{Contraindications and Precautions}
The team reviews nasal obstruction, severe swallowing difficulty, recent nasal or esophageal injury, bleeding risk, and the ability to cooperate with swallowing instructions. Significant infection or instability may require postponement or another test.

\section*{When to Seek Medical Care}
Follow the procedure team's specific instructions. Seek urgent medical assessment for severe or worsening pain, uncontrolled bleeding, fever or signs of infection, trouble breathing, chest pain, fainting, new weakness or confusion, inability to urinate, or any rapidly worsening symptom. The appropriate warning signs depend on the procedure and the patient's medical condition.

\section*{References}
\begin{enumerate}
  \item MedlinePlus. Esophageal manometry. U.S. National Library of Medicine; 2025.
  \item Current Medical Diagnosis and Treatment. McGraw-Hill; 2025.
\end{enumerate}
\clearpage
